\begin{landscape}
\begingroup\renewcommand{\baselinestretch}{1}\footnotesize\setlength{\LTcapwidth}{\linewidth}

\begin{ThreePartTable}
\begin{TableNotes}
\item Note: Data are from the American Community Survey (ACS), 2011 to 2013. The sample is adults aged 26 to 64 living in households, with family income at or below 138\% of the federal poverty level. It excludes adults with Medicare, US citizens born in the US territories or abroad to American parents, non-citizens in the US fewer than 5 years, and Delaware, the District of Columbia, Massachusetts, New York and Vermont, which covered low-income adults before 2014. Values are percentages, except where marked mean (SD), weighted with ACS person weights; n is the unweighted number of respondents, and the number of states is in parentheses. Any insurance follows the Census Bureau definition, which does not count Indian Health Service coverage alone. Immigrant characteristics are shown for the foreign-born only. State context values are person-weighted means of the state's annual unemployment rate (BLS), its Immigrant Policy Climate index without the health care items (Samari), and the average ADA score (0 to 100) of its members of Congress. 
\item Abbreviations: ACS, American Community Survey; ADA, Americans for Democratic Action; BLS, Bureau of Labor Statistics; SD, standard deviation; US, United States.
\end{TableNotes}
\begin{longtable}[t]{lrrrrrr}
\caption{Characteristics of low-income adults before Medicaid expansion by citizenship and state expansion status}\\
\toprule
\multicolumn{1}{c}{ } & \multicolumn{3}{c}{Non-expansion states (17)} & \multicolumn{3}{c}{Expansion states (29)} \\
\cmidrule(l{3pt}r{3pt}){2-4} \cmidrule(l{3pt}r{3pt}){5-7}
 & \makecell[br]{US-born\\n = 213,308} & \makecell[br]{Naturalized\\citizen\\n = 14,811} & \makecell[br]{Non-citizen\\n = 35,054} & \makecell[br]{US-born\\n = 282,783} & \makecell[br]{Naturalized\\citizen\\n = 29,071} & \makecell[br]{Non-citizen\\n = 58,319}\\
\midrule
\endfirsthead
\caption[]{Characteristics of low-income adults before Medicaid expansion by citizenship and state expansion status \textit{(continued)}}\\
\toprule
 & \makecell[br]{US-born\\n = 213,308} & \makecell[br]{Naturalized\\citizen\\n = 14,811} & \makecell[br]{Non-citizen\\n = 35,054} & \makecell[br]{US-born\\n = 282,783} & \makecell[br]{Naturalized\\citizen\\n = 29,071} & \makecell[br]{Non-citizen\\n = 58,319}\\
\midrule
\endhead

\endfoot
\bottomrule
\insertTableNotes
\endlastfoot
\addlinespace[0.3em]
\multicolumn{7}{l}{\textbf{Health insurance coverage}}\\
\hspace{1em}Any insurance & 52.7 & 44.1 & 20.7 & 61.4 & 60.9 & 33.4\\
\hspace{1em}Medicaid & 25.3 & 17.2 & 8.3 & 34.5 & 33.1 & 19.6\\
\hspace{1em}Employer-sponsored insurance & 20.3 & 20.5 & 10.1 & 19.8 & 21.5 & 11.9\\
\hspace{1em}Direct purchase & 7.0 & 7.2 & 2.7 & 7.9 & 8.2 & 2.8\\
\addlinespace[0.3em]
\multicolumn{7}{l}{\textbf{Demographics}}\\
\hspace{1em}Age, mean (SD) & 42.6 (11.2) & 44.8 (10.3) & 39.8 (9.4) & 42.6 (11.2) & 45.5 (10.4) & 40.1 (9.4)\\
\hspace{1em}Female & 58.2 & 56.8 & 52.5 & 57.0 & 56.6 & 53.6\\
\hspace{1em}Race and ethnicity &  &  &  &  &  & \\
\hspace{1em}\hspace{1em}White & 58.3 & 11.4 & 3.9 & 64.2 & 18.0 & 5.5\\
\hspace{1em}\hspace{1em}Black & 28.4 & 14.3 & 5.2 & 19.6 & 7.4 & 2.8\\
\hspace{1em}\hspace{1em}Asian & 0.2 & 16.4 & 4.1 & 0.8 & 29.2 & 8.2\\
\hspace{1em}\hspace{1em}Hispanic & 10.0 & 55.4 & 86.0 & 11.0 & 42.9 & 82.4\\
\hspace{1em}\hspace{1em}Other & 3.2 & 2.6 & 0.8 & 4.4 & 2.5 & 1.1\\
\hspace{1em}Married & 32.0 & 57.1 & 59.4 & 28.7 & 58.1 & 56.2\\
\hspace{1em}Disability & 20.9 & 10.5 & 6.4 & 21.2 & 11.5 & 6.7\\
\addlinespace[0.3em]
\multicolumn{7}{l}{\textbf{Socioeconomic status}}\\
\hspace{1em}Education &  &  &  &  &  & \\
\hspace{1em}\hspace{1em}Less than high school & 21.4 & 29.4 & 59.0 & 18.6 & 32.9 & 59.8\\
\hspace{1em}\hspace{1em}High school & 36.4 & 29.2 & 24.4 & 36.0 & 26.3 & 23.4\\
\hspace{1em}\hspace{1em}Some college or Associate degree & 32.3 & 25.0 & 10.4 & 33.7 & 24.1 & 10.6\\
\hspace{1em}\hspace{1em}College degree & 7.6 & 11.9 & 4.5 & 8.7 & 12.0 & 4.4\\
\hspace{1em}\hspace{1em}Graduate and beyond & 2.3 & 4.6 & 1.6 & 3.0 & 4.7 & 1.8\\
\hspace{1em}Employment &  &  &  &  &  & \\
\hspace{1em}\hspace{1em}Employed & 44.4 & 56.6 & 59.1 & 41.7 & 52.8 & 55.8\\
\hspace{1em}\hspace{1em}Unemployed & 15.9 & 11.9 & 9.1 & 16.9 & 13.0 & 10.8\\
\hspace{1em}\hspace{1em}Not in labor force & 39.7 & 31.4 & 31.8 & 41.4 & 34.2 & 33.4\\
\hspace{1em}Income as \% of federal poverty level &  &  &  &  &  & \\
\hspace{1em}\hspace{1em}0-100\% & 67.9 & 60.3 & 65.0 & 69.4 & 60.6 & 64.2\\
\hspace{1em}\hspace{1em}101-138\% & 32.1 & 39.7 & 35.0 & 30.6 & 39.4 & 35.8\\
\addlinespace[0.3em]
\multicolumn{7}{l}{\textbf{Immigrant characteristics}}\\
\hspace{1em}Years in the US, mean (SD) &  & 23.4 (11.3) & 15.7 (8.4) &  & 23.9 (10.9) & 16.8 (8.8)\\
\hspace{1em}Lived $>$25\% of life in the US &  & 87.5 & 75.1 &  & 88.8 & 78.4\\
\hspace{1em}English proficiency &  &  &  &  &  & \\
\hspace{1em}\hspace{1em}Only English &  & 14.9 & 5.8 &  & 11.3 & 5.2\\
\hspace{1em}\hspace{1em}Very well &  & 35.6 & 14.8 &  & 32.2 & 14.6\\
\hspace{1em}\hspace{1em}Well &  & 26.6 & 20.3 &  & 28.7 & 20.6\\
\hspace{1em}\hspace{1em}Not well &  & 18.8 & 37.0 &  & 23.1 & 37.8\\
\hspace{1em}\hspace{1em}Not at all &  & 4.0 & 22.1 &  & 4.7 & 21.8\\
\hspace{1em}Region of birth &  &  &  &  &  & \\
\hspace{1em}\hspace{1em}Mexico &  & 29.0 & 63.9 &  & 30.0 & 66.5\\
\hspace{1em}\hspace{1em}Central America &  & 8.1 & 13.4 &  & 7.2 & 12.3\\
\hspace{1em}\hspace{1em}Caribbean &  & 21.6 & 9.1 &  & 4.6 & 2.2\\
\hspace{1em}\hspace{1em}South America &  & 8.9 & 4.6 &  & 3.9 & 3.0\\
\hspace{1em}\hspace{1em}East and Southeast Asia &  & 12.6 & 3.1 &  & 25.6 & 7.3\\
\hspace{1em}\hspace{1em}South and Central Asia &  & 4.2 & 1.2 &  & 5.7 & 1.4\\
\hspace{1em}\hspace{1em}Middle East and North Africa &  & 4.8 & 0.7 &  & 8.5 & 1.6\\
\hspace{1em}\hspace{1em}Sub-Saharan Africa &  & 4.1 & 1.5 &  & 5.1 & 2.0\\
\hspace{1em}\hspace{1em}Europe, Canada and Oceania &  & 6.7 & 2.5 &  & 9.4 & 3.8\\
\addlinespace[0.3em]
\multicolumn{7}{l}{\textbf{State context}}\\
\hspace{1em}State unemployment rate, mean (SD) & 7.9 (1.5) & 7.9 (1.4) & 7.8 (1.5) & 8.3 (1.6) & 9.3 (1.7) & 9.4 (1.7)\\
\hspace{1em}Immigrant Policy Climate, mean (SD) & -2.7 (1.5) & -2.5 (1.3) & -2.2 (1.4) & -0.8 (2.1) & 0.4 (2.1) & 0.5 (2.2)\\
\hspace{1em}ADA score, mean (SD) & 23.7 (9.5) & 25.0 (6.7) & 24.3 (7.3) & 44.3 (17.6) & 53.3 (13.9) & 53.2 (13.7)\\*
\end{longtable}
\end{ThreePartTable}
\endgroup
\end{landscape}
